\section{Conclusion}

We began with an empirical observation: when ensembling iterative reasoners,
selecting by \emph{completion speed} dramatically outperforms probability
averaging. This led us to a hypothesis---inference speed is an implicit
confidence signal---and a method: winner-take-all training that internalizes
ensemble diversity within a single model.

% Source: x182a.log (best seed 97.64%), x182{a,b,c,d} mean 96.9% ± 0.6%
% Baseline x182{m,n,o,p}: 85.5% ± 1.3%
Our results on Sudoku-Extreme are striking: a single 7M parameter model achieves
96.9\% $\pm$ 0.6\% puzzle accuracy (best seed: 97.6\%), matching a 3-model
halt-based ensemble and improving 11pp over the TRM baseline. The key insight is
not architectural but algorithmic: by letting $K{=}4$ parallel hypotheses
compete during training, we teach the model to find confident solutions faster.

The broader implication is that \textbf{how quickly} a model converges may be
as informative as \textbf{what} it outputs. This connects to a rich literature
on speed-accuracy tradeoffs in human cognition, where fast decisions correlate
with confidence. Adaptive computation mechanisms like ACT may be learning to
detect this same signal.

\paragraph{Future directions.}
\begin{itemize}
  \item \textbf{Beyond Sudoku.} Validating on ARC, maze-solving, and other
    reasoning benchmarks.
  \item \textbf{Scaling $K$.} Understanding why $K{=}4$ saturates and whether
    harder domains require more heads.
  \item \textbf{Head specialization.} Encouraging heads to learn diverse
    strategies rather than converging to similar solutions.
  \item \textbf{Theoretical foundations.} Formalizing the connection between
    convergence speed and solution quality.
\end{itemize}
